\subsection{\texorpdfstring{$\Xi$}{Xi}/Hardy 实值化、dyadic 首壳与后验实部压缩的刚性链}\label{subsec:conclusion-xi-hardy-dyadic-shell-posterior-sigma-rigidity}
沿用小节 \ref{subsec:conclusion-xi-dyadic-recursion-zero-stability}、小节 \ref{subsec:conclusion-phase-prefix-mellin-transverse-gabcke-rigidity}、\S\ref{subsec:cdim-rs-gabcke-zero-recursion}、\S\ref{subsec:cdim-afe-posterior-sigma-contraction} 与小节 \ref{subsubsec:xi-riemann-siegel-recursive-zero-transfer} 的记号。为把完成 \(\zeta\) 核的 dyadic 余项预算与 Hardy 实值化侧的 Riemann--Siegel--Gabcke 零点递归并置，下文把两条层级都以共同的分辨率索引 \(K\) 记之：在完成核侧，\(R_K(s)\) 表示 dyadic 二进壳余项；在 Hardy 侧，\(Z^{(K)}(t)\) 表示第 \(K\) 阶 Riemann--Siegel--Gabcke 截断，\(\tau_n^{(K)}\) 表示其第 \(n\) 个正零点。再记
\[
E_K(T):=\sup_{|t|\le T}\left|R_K\!\left(\frac12+\mathrm{i}t\right)\right|.
\]
前述正文已分别给出：完成核 dyadic 余项的不完全伽马展开与首项相位律、Rouch\'e--Taylor 零点传输、层级零点的单侧收敛猜想，以及对称近似函数方程诱导的后验实部约束与收缩吸引子判据。把这些接口交叉压缩后，可把 \(\Xi\)/Hardy 模块进一步收束为一条独立的分辨率几何链：dyadic 首壳决定余项主相位；零点塔在可控预算下呈超指数 Cauchy 尾；而后验实部偏移则只通过“截断比率 + \(\chi\)-项余差”这两个标量被压缩，并相较横向零漂移额外损失一个 \(\log t\) 因子。

\begin{theorem}[dyadic 余项壳的首项包络律]\label{thm:conclusion-xi-dyadic-shell-envelope-law}
对任意固定 \(T>0\)，当 \(K\to\infty\) 时，在 \(|t|\le T\) 上一致成立
\[
R_K\!\left(\frac12+\mathrm{i}t\right)
=
\frac{2}{\pi}\,2^{-3K/4}e^{-\pi 2^K}
\cos\!\left(\frac{K\log 2}{2}\,t\right)
+ O_T\!\left(2^{-7K/4}e^{-\pi 2^K}\right)
+ O\!\left(2^{-3K/4}e^{-4\pi 2^K}\right).
\]
等价地，
\[
R_K\!\left(\frac12+\mathrm{i}t\right)
=
\frac{2}{\pi}\,2^{-3K/4}e^{-\pi 2^K}
\cos\!\left(\frac{K\log 2}{2}\,t\right)
+ \Bigl[O_T(2^{-K})+O(e^{-3\pi 2^K})\Bigr]\,
2^{-3K/4}e^{-\pi 2^K}.
\]
特别地，
\[
E_K(T)
=
\frac{2}{\pi}\,2^{-3K/4}e^{-\pi 2^K}
\Bigl(1+O_T(2^{-K})+O(e^{-3\pi 2^K})\Bigr).
\]
\end{theorem}
\begin{proof}
不完全伽马展开与临界线首项相位律正是定理 \ref{thm:xi-completed-zeta-remainder-phase-law} 的内容。对 \(|t|\le T\) 取上确界时，主项振荡因子
\(
\cos\bigl(\frac{K\log2}{2}t\bigr)
\)
在 \(t=0\) 处达到绝对值 \(1\)，从而得到所示包络公式。证毕。
\end{proof}

\begin{corollary}[每增一层的二进制精度收益律]\label{cor:conclusion-xi-dyadic-shell-bit-gain-law}
\leanverified{paper\_conclusion\_xi\_dyadic\_shell\_bit\_gain\_law}
记
\[
B_K(T):=-\log_2 E_K(T).
\]
则
\[
B_{K+1}(T)-B_K(T)
=
\frac{\pi}{\log 2}\,2^K+\frac34+o(1)
\qquad (K\to\infty).
\]
\end{corollary}
\begin{proof}
由定理 \ref{thm:conclusion-xi-dyadic-shell-envelope-law}，
\[
\frac{E_{K+1}(T)}{E_K(T)}
=
2^{-3/4}e^{-\pi 2^K}\bigl(1+o(1)\bigr).
\]
两边取 \(-\log_2\) 即得结论。证毕。
\end{proof}

\begin{theorem}[dyadic 首壳的零点晶格]\label{thm:conclusion-xi-dyadic-shell-zero-lattice}
\leanverified{paper\_conclusion\_xi\_dyadic\_shell\_zero\_lattice}
记主导余项壳
\[
\mathcal R_K(t):=
\frac{2}{\pi}\,2^{-3K/4}e^{-\pi 2^K}
\cos\!\left(\frac{K\log 2}{2}\,t\right).
\]
则其零点集合为
\[
t_{K,\ell}=\frac{(2\ell+1)\pi}{K\log 2},
\qquad \ell\in\ZZ,
\]
并且相邻零点间距恰为
\[
\Delta t_K=\frac{2\pi}{K\log 2}.
\]
\end{theorem}
\begin{proof}
直接解
\[
\cos\!\left(\frac{K\log2}{2}t\right)=0
\]
即可。证毕。
\end{proof}

\begin{theorem}[相邻截断零点的 Rouch\'e--Taylor 传输不等式]\label{thm:conclusion-xi-adjacent-truncation-rouche-taylor-transport}
设 \(\tau_n^{(K)}\) 是 \(Z^{(K)}\) 的一个单零点。若在闭圆盘 \(\overline{D(\tau_n^{(K)},r_K)}\) 上定义
\[
\varepsilon_K:=
\sup_{|t-\tau_n^{(K)}|\le r_K}\bigl|Z^{(K+1)}(t)-Z^{(K)}(t)\bigr|,
\qquad
M_K:=
\sup_{|t-\tau_n^{(K)}|\le r_K}\bigl|(Z^{(K)})''(t)\bigr|,
\]
并满足
\[
\bigl|(Z^{(K)})'(\tau_n^{(K)})\bigr|>\frac{M_K}{2}r_K,
\qquad
\varepsilon_K<
\left(\bigl|(Z^{(K)})'(\tau_n^{(K)})\bigr|-\frac{M_K}{2}r_K\right)r_K,
\]
则 \(Z^{(K+1)}\) 在该圆盘内恰有一个单零点 \(\tau_n^{(K+1)}\)，且
\[
\bigl|\tau_n^{(K+1)}-\tau_n^{(K)}\bigr|
\le
\frac{\varepsilon_K}{
\bigl|(Z^{(K)})'(\tau_n^{(K)})\bigr|-\frac{M_K}{2}r_K }.
\]
\end{theorem}
\begin{proof}
对定理 \ref{thm:xi-rouche-taylor-zero-transfer} 取
\[
g=Z^{(K)},\qquad f=Z^{(K+1)},\qquad t_0=\tau_n^{(K)}
\]
即得。证毕。
\end{proof}

\begin{corollary}[零点塔的超指数 Cauchy 尾]\label{cor:conclusion-xi-zero-tower-superexp-cauchy-tail}
\leanverified{paper\_conclusion\_xi\_zero\_tower\_superexp\_cauchy\_tail}
若对某一固定零点序号 \(n\) 存在常数 \(c_T>0\)、\(C_T>0\) 与层级下界 \(K_0\)，使得对所有 \(K\ge K_0\)：
\[
\bigl|(Z^{(K)})'(\tau_n^{(K)})\bigr|-\frac{M_K}{2}r_K\ge c_T,
\qquad
\varepsilon_K\le C_T E_K(T),
\]
则
\[
\sum_{j\ge K}\bigl|\tau_n^{(j+1)}-\tau_n^{(j)}\bigr|<\infty,
\]
并且
\[
\sum_{j\ge K}\bigl|\tau_n^{(j+1)}-\tau_n^{(j)}\bigr|
=
O_T\!\bigl(E_K(T)\bigr).
\]
若再假设 \(\tau_n^{(K)}\to \gamma_n\)，则
\[
\bigl|\gamma_n-\tau_n^{(K)}\bigr|
=
O_T\!\bigl(E_K(T)\bigr).
\]
\end{corollary}
\begin{proof}
由定理 \ref{thm:conclusion-xi-adjacent-truncation-rouche-taylor-transport}，
\[
\bigl|\tau_n^{(K+1)}-\tau_n^{(K)}\bigr|
\le c_T^{-1}C_T E_K(T).
\]
而由推论 \ref{cor:conclusion-xi-dyadic-shell-bit-gain-law}，\(E_{K+1}(T)/E_K(T)=2^{-3/4}e^{-\pi 2^K}(1+o(1))\)，故 \(\sum_{j\ge K}E_j(T)\) 被首项 \(E_K(T)\) 支配，从而尾和为 \(O_T(E_K(T))\)。若再有 \(\tau_n^{(K)}\to\gamma_n\)，则
\[
\bigl|\gamma_n-\tau_n^{(K)}\bigr|
\le
\sum_{j\ge K}\bigl|\tau_n^{(j+1)}-\tau_n^{(j)}\bigr|
=
O_T\!\bigl(E_K(T)\bigr).
\]
证毕。
\end{proof}

\begin{theorem}[后验实部偏移的双标量压缩律]\label{thm:conclusion-xi-posterior-realpart-two-scalar-compression}
设 \(s=\sigma+\mathrm{i}t\) 为 \(\zeta\) 的非平凡零点，\(t\ge 2\)。在定理 \ref{thm:cdim-afe-posterior-sigma-constraint} 的余项证书条件下，记
\[
N:=\Bigl\lfloor \sqrt{\frac{t}{2\pi}}\Bigr\rfloor,
\qquad
\eta(s):=\frac{\mathcal E(s)}{|A_N(s)|}\in(0,1),
\]
以及
\[
\Delta_\chi(t,\sigma):=
\log|\chi(\sigma+\mathrm{i}t)|
-\Bigl(\frac12-\sigma\Bigr)\log\frac{t}{2\pi}.
\]
若存在区间 \(J\subset(0,1)\) 使 \(\sigma\in J\)、\(\Phi_t\) 在 \(J\) 上可微且
\[
\sup_{\xi\in J}|\Phi_t'(\xi)|\le q<1,
\]
则
\[
\left|\sigma-\frac12\right|
\le
\frac{1}{1-q}
\left[
\frac{1}{\log(t/2\pi)}
\log\frac{1+\eta(s)}{1-\eta(s)}
+
\frac{|\Delta_\chi(t,\sigma)|}{\log(t/2\pi)}
\right].
\]
\end{theorem}
\begin{proof}
这正是定理 \ref{thm:cdim-afe-contraction-attractor} 对其误差项
\[
\varepsilon(t,\sigma)=
\frac{1}{\log(t/2\pi)}
\log\!\Bigl(\frac{1+\eta(s)}{1-\eta(s)}\Bigr)
+
\frac{|\Delta_\chi(t,\sigma)|}{\log(t/2\pi)}
\]
的直接展开。证毕。
\end{proof}

\begin{corollary}[小余项制度下的线性化偏移界]\label{cor:conclusion-xi-small-remainder-linearized-sigma-bound}
若 \(0<\eta(s)\le \eta_0<1\)，则
\[
\log\frac{1+\eta(s)}{1-\eta(s)}
=
2\eta(s)+O_{\eta_0}\!\bigl(\eta(s)^3\bigr),
\]
从而
\[
\left|\sigma-\frac12\right|
\le
\frac{
2\eta(s)+O_{\eta_0}\!\bigl(\eta(s)^3\bigr)+|\Delta_\chi(t,\sigma)|
}{
(1-q)\log(t/2\pi)
}.
\]
特别地，若 \(\eta(s)\to0\) 且 \(\Delta_\chi(t,\sigma)=O(1/t)\)，则
\[
\left|\sigma-\frac12\right|
=
O\!\left(\frac{\eta(s)}{\log t}\right)+O\!\left(\frac{1}{t\log t}\right).
\]
\end{corollary}
\begin{proof}
对奇函数
\[
\log\frac{1+\eta}{1-\eta}
\]
在 \(\eta=0\) 处作 Taylor 展开即可，再代入定理 \ref{thm:conclusion-xi-posterior-realpart-two-scalar-compression}。证毕。
\end{proof}

\begin{theorem}[dyadic 余项预算导出的纵向超指数收缩]\label{thm:conclusion-xi-dyadic-budget-vertical-superexp-contraction}
在定理 \ref{thm:conclusion-xi-posterior-realpart-two-scalar-compression} 的设定下，若在零点窗口 \(2\le t\le T\) 内有一致下界
\[
|A_N(s)|\ge a_0>0,
\]
并且余项证书满足
\[
\mathcal E(s)\le E_K(T),
\]
则
\[
\eta(s)=O_T\!\bigl(2^{-3K/4}e^{-\pi 2^K}\bigr).
\]
进一步，若 \(\Delta_\chi(t,\sigma)=O(1/t)\)，则
\[
\left|\sigma-\frac12\right|
\le
\frac{C_T}{(1-q)\log(t/2\pi)}\,2^{-3K/4}e^{-\pi 2^K}
+
O\!\left(\frac{1}{t\log t}\right).
\]
\end{theorem}
\begin{proof}
由定理 \ref{thm:conclusion-xi-dyadic-shell-envelope-law}，
\[
\mathcal E(s)\le E_K(T)=O_T\!\bigl(2^{-3K/4}e^{-\pi 2^K}\bigr).
\]
再除以 \(|A_N(s)|\ge a_0\) 即得 \(\eta(s)\) 的同阶上界。将其代入推论 \ref{cor:conclusion-xi-small-remainder-linearized-sigma-bound}，并用 \(\Delta_\chi(t,\sigma)=O(1/t)\) 吸收余项，即得结论。证毕。
\end{proof}

\begin{corollary}[零点认证的比特复杂度只呈对数增长]\label{cor:conclusion-xi-zero-certification-logarithmic-depth}
在推论 \ref{cor:conclusion-xi-zero-tower-superexp-cauchy-tail} 的设定下，若要求
\[
|\gamma_n-\tau_n^{(K)}|\le 2^{-B},
\]
则充分大的 \(K\) 只需满足
\[
K=\log_2 B-\log_2\!\left(\frac{\pi}{\log 2}\right)+o(1)
\qquad (B\to\infty).
\]
换言之，dyadic 层级深度对所需二进制精度只呈对数增长。
\end{corollary}
\begin{proof}
由推论 \ref{cor:conclusion-xi-zero-tower-superexp-cauchy-tail}，
\[
|\gamma_n-\tau_n^{(K)}|=O_T(E_K(T)).
\]
而由定理 \ref{thm:conclusion-xi-dyadic-shell-envelope-law}，
\[
-\log_2 E_K(T)
=
\frac{\pi}{\log2}\,2^K+\frac34K+O(1).
\]
令右端与 \(B\) 同阶平衡，便得到所示 \(K\)-尺度。证毕。
\end{proof}

\begin{theorem}[横向零漂移与纵向实部漂移的 \texorpdfstring{$\log t$}{log t}-级分裂]\label{thm:conclusion-xi-horizontal-vertical-logt-splitting}
在推论 \ref{cor:conclusion-xi-zero-tower-superexp-cauchy-tail} 与定理 \ref{thm:conclusion-xi-dyadic-budget-vertical-superexp-contraction} 的共同假设下，若再存在 \(c_T>0\) 使
\[
|\gamma_n-\tau_n^{(K)}|\ge c_T E_K(T),
\]
则
\[
\frac{\left|\sigma-\frac12\right|}{|\gamma_n-\tau_n^{(K)}|}
\le
O\!\left(\frac{1}{\log t}\right)
+
O\!\left(\frac{1}{t\log t}\,E_K(T)^{-1}\right).
\]
特别地，在固定 \(K\) 且 \(t\to\infty\) 的 regime 中，纵向偏移比横向零漂移严格低一个 \(\log t\) 因子。
\end{theorem}
\begin{proof}
由定理 \ref{thm:conclusion-xi-dyadic-budget-vertical-superexp-contraction}，
\[
\left|\sigma-\frac12\right|
=
O\!\left(\frac{E_K(T)}{\log t}\right)
+
O\!\left(\frac{1}{t\log t}\right).
\]
再用假设 \( |\gamma_n-\tau_n^{(K)}|\ge c_T E_K(T)\) 即得
\[
\frac{\left|\sigma-\frac12\right|}{|\gamma_n-\tau_n^{(K)}|}
\le
O\!\left(\frac{1}{\log t}\right)
+
O\!\left(\frac{1}{t\log t}\,E_K(T)^{-1}\right).
\]
证毕。
\end{proof}

\begin{theorem}[在单侧收敛与首省略项方向猜想下的壳相位定向律]\label{thm:conclusion-xi-shell-phase-direction-under-conjectures}
假设猜想 \ref{conj:cdim-rs-gabcke-one-sided-convergence} 与猜想 \ref{conj:xi-riemann-siegel-omitted-term-direction-law} 同时成立。则对每个固定零点序号 \(n\)，存在 \(K_0(n)\) 使得对所有 \(K\ge K_0(n)\)，只要首壳主项严格支配次级误差，即
\[
\left|
\cos\!\left(\frac{K\log2}{2}\,\tau_n^{(K)}\right)
\right|
>
C_T\bigl(2^{-K}+e^{-3\pi 2^K}\bigr),
\]
就有
\[
\operatorname{sgn}\bigl(\tau_n^{(K)}-\gamma_n\bigr)
=
-\operatorname{sgn}\!\Bigl(\Delta_K(\tau_n^{(K)})\Bigr)
=
-\operatorname{sgn}\!\left(
\cos\!\left(\frac{K\log2}{2}\,\tau_n^{(K)}\right)
\right).
\]
因此，在高层 regime 中，零点位移方向由 dyadic 首壳的局部相位单独锁定。
\end{theorem}
\begin{proof}
由猜想 \ref{conj:xi-riemann-siegel-omitted-term-direction-law}，
\[
\operatorname{sgn}\bigl(\tau_n^{(K)}-\gamma_n\bigr)
=
-\operatorname{sgn}\!\Bigl(\Delta_K(\tau_n^{(K)})\Bigr).
\]
另一方面，定理 \ref{thm:xi-completed-zeta-remainder-phase-law} 给出 dyadic 余项的首壳展开：在固定窗口内，
\[
\Delta_K(\tau_n^{(K)})
=
\frac{2}{\pi}\,2^{-3K/4}e^{-\pi 2^K}
\cos\!\left(\frac{K\log2}{2}\,\tau_n^{(K)}\right)
+ \Bigl[O_T(2^{-K})+O(e^{-3\pi 2^K})\Bigr]\,
2^{-3K/4}e^{-\pi 2^K}.
\]
若主余弦项严格支配括号中的次级修正，则 \(\Delta_K(\tau_n^{(K)})\) 与该余弦项同号，于是得到所示结论。单侧收敛假设保证这一定向律与高层同侧逼近方向相容。证毕。
\end{proof}

\begin{conjecture}[dyadic 壳相位刚性]\label{conj:conclusion-xi-dyadic-shell-phase-rigidity}
对每个充分高的临界线零点 \(\gamma_n\)，存在 \(K_0(n)\) 使得对所有 \(K\ge K_0(n)\)，都有严格相位锁定
\[
\operatorname{sgn}\bigl(\tau_n^{(K)}-\gamma_n\bigr)
=
-\operatorname{sgn}\!\left(
\cos\!\left(\frac{K\log2}{2}\,\gamma_n\right)
\right).
\]
亦即：一旦进入高层 regime，第 \(K\) 层零点位移的方向不再由全部余项共同决定，而由 dyadic 首壳在真零点处的相位单独决定。
\end{conjecture}

\paragraph{统一读法}
这一组结果把 \(\Xi\)/Hardy/Riemann--Siegel 模块中的三条原始接口压成了一条单独的分辨率几何链。首先，完成核 dyadic 余项并非普通截断尾项，而是由单一首壳主导的超指数预算对象；其主相位由
\(
\cos(\frac{K\log2}{2}t)
\)
锁定，并在每推进一层时带来 \(\asymp 2^K\) 的二进制精度收益。其次，Hardy 侧的截断零点并不是松散数值近似，而是在 Rouch\'e--Taylor 传输下形成具有超指数 Cauchy 尾的零点塔；若相邻层漂移确实饱和该预算，则零点认证所需层级深度只随目标位数对数增长。最后，对称近似函数方程把离临界线方向的实部偏移压缩成“截断比率 \(\eta\) + \(\chi\)-项余差 \(\Delta_\chi\)”这两个标量，而在 dyadic 预算下该纵向偏移又比横向零漂移额外低一个 \(\log t\) 因子。故在这一模块里，dyadic 层级并不是普通截断，而是一种同时携带首壳主导、超指数尾、零点塔稳定性与后验实部压缩的可审计分辨率几何。

\endinput
